\section{Classical pencils and the small affine exclusions}
\begin{proposition}[Classical local actions]\label{classical}
The natural groups containing $\operatorname{PSL}_{m+1}(q)$ on
$\operatorname{PG}(m,q)$, or containing translations and
$\operatorname{SL}_m(q)$ on $\operatorname{AG}(m,q)$, are flag-transitive
and locally two-transitive for $m\ge2$. In the affine plane the induced
special-linear pencil group is $\operatorname{PSL}_2(q)$; it need not
contain $\operatorname{PGL}_2(q)$.
\end{proposition}
\begin{proof}
Special-linear matrices are transitive on nonzero vectors and on ordered
pairs of distinct one-dimensional subspaces: choose independent
representatives, extend them to bases, and rescale a representative or a
complementary vector to correct the determinant. In the projective pencil
at a point, any quotient matrix $A$ extends to the determinant-one matrix
$\operatorname{diag}((\det A)^{-1},A)$. On a projective line, any
$\operatorname{GL}_2(q)$ action extends by its inverse determinant on a
complementary coordinate. Both local projective actions are therefore
two-transitive. The affine pencil is precisely the projective action
on vector lines. Its vector-line stabilizer contains every nonzero
multiplier on that line, compensating the determinant on a complement;
line translations give two-transitivity there. Supergroups preserve these
properties.

For $m=2$ the scalar kernel in $\operatorname{SL}_2(q)$ has order
$\gcd(2,q-1)$, giving exactly $\operatorname{PSL}_2(q)$ on the pencil.
For odd $q$ this has index two in $\operatorname{PGL}_2(q)$.
Thus $\operatorname{ASL}_2(3)$ on $\operatorname{AG}(2,3)$ has pencil
group $A_4$, not $S_4$, despite being locally two-transitive. The
$\operatorname{PGL}_2(q)$ inclusions in source Propositions 3.7 and 3.8
must be replaced by $\operatorname{PSL}_2(q)$. The source's final
special-linear branches do not require this erroneous inclusion.
\end{proof}

\begin{proposition}[Natural symplectic pencils]\label{symplectic}
For even $m\ge4$, the natural $\operatorname{Sp}_m(q)$ is primitive but
not two-transitive on vector lines. Its normalizer over $\F_p$ preserves
perpendicularity, up to field automorphism and nonzero form multiplier.
Consequently every group between it and its normalizer has a locally
primitive, but not locally two-transitive, affine-space action.
\end{proposition}
\begin{proof}
Symplectic basis extension shows that the stabilizer of a vector line
$P$ has three orbits on vector lines: $P$, the remaining points of
$P^\perp$, and those outside $P^\perp$. Write
\[
 v=(q^m-1)/(q-1),\qquad a=(q^{m-1}-1)/(q-1).
\]
A block through $P$ is a union of stabilizer orbits, so its possible
proper nontrivial sizes are $a$ and $q^{m-1}+1$. But
$\gcd(a,v)=1$, and $v/2<q^{m-1}+1<v$; neither size divides $v$.
This proves primitivity, and the two nonempty other orbits disprove
two-transitivity.

The commuting algebra over $\F_p$ is $\F_q$. Indeed commuting with
every transvection $x\mapsto x+t\langle x,w\rangle w$ forces an
endomorphism to preserve each $\F_qw$, the image after subtracting the
identity. An additive map preserving all vector lines is a common scalar:
apply it to $u,v,u+v,u+cv$ for independent $u,v$ and $c\in\F_q$.
Lemma~\ref{field} now makes a normalizer element semilinear over $\F_q$.
After removing its field automorphism, pulling back the alternating form
gives another invariant form. Relative to the original nondegenerate
form its defining operator commutes with the group, and hence is a
scalar. Perpendicularity is preserved even by the full normalizer.
Finally the symplectic group is transitive on nonzero vectors, so a
vector-line stabilizer is transitive on the nonzero vectors of that
line. Proposition~\ref{translation} gives the remaining local assertions.
\end{proof}

\begin{lemma}[A small normal subgroup test]\label{normaltest}
A nontrivial normal subgroup in a faithful primitive permutation group
is transitive. In particular its order is at least the degree.
A faithful primitive group with a nontrivial center is regular of prime
order. These tests apply to the induced group after removal of the
restriction kernel.
\end{lemma}
\begin{proof}
Normal-subgroup orbits are a block system. All singleton orbits would
make the subgroup trivial by faithfulness. A central transitive subgroup
is regular, since an element fixing one point fixes all. Its centralizer
in the symmetric group is its own regular copy, forcing the whole group
to be abelian regular. Primitivity of a regular group is equivalent to
having no proper nontrivial subgroup, hence to prime order.
\end{proof}

Retain Huybrechts' actual representation cases as quoted in source
Proposition 3.4, rather than just abstract isomorphic subgroup names.
The exceptional degree-364 pencil group is $\operatorname{PSL}_2(13)$,
of order $1092$. A point stabilizer has order three and lies properly
in a cyclic split torus of order six: an order-three element diagonalizes
over $\F_{13}$, and the determinant-one diagonal torus modulo scalars
has order six. Thus this action is imprimitive. In the degree-40 cases
$2^4.A$, the normal subgroup of order sixteen has a nontrivial image
of order at most sixteen on directions: the scalar kernel in the
linear $\F_3$-action has order at most two. Such an image cannot be
transitive of degree forty, so Lemma~\ref{normaltest} excludes
primitivity. In the remaining $4.A_5$ or $N.S_5$ cases ($|N|\leq4$),
the specified normal subgroup of order at most four likewise has either
a nontrivial, nontransitive image, or lies entirely in the restriction
kernel. In the latter case the induced group is a quotient of $A_5$
or $S_5$. The former order is not divisible by forty. A transitive
$S_5$ quotient of degree forty must be $S_5$ itself (its other proper
quotients have order at most two); its order-three point stabilizer
lies properly in an $S_3$. Thus this action is imprimitive as well.
This supplies the missing explicit local exclusions, conditional on
the retained representation classification.

\begin{proposition}[The exceptional affine plane]\label{nineplane}
Retain Foulser's ten flag-transitive actions on $\operatorname{AG}(2,9)$
with induced pencil group $A_5$ or $S_5$, each containing the specified
minimal $F=3^4\rtimes2.A_5$ in its natural $\operatorname{SL}_2(9)$
representation. All ten are locally primitive and none is locally
two-transitive.
\end{proposition}
\begin{proof}
In the degree-ten action $A_5$ has point stabilizer $S_3$, which is
maximal. Indeed a proper overgroup could only have order twelve or
thirty; an order-twelve subgroup is $A_4$ and contains no $S_3$, and
simplicity forbids index two. Neither $60$ nor $120$ is divisible by
$10\cdot9$, excluding two-transitivity for either pencil image.

Choose an involution of $A_5$ and a lift $f$ in $2.A_5$. The only
involution in $\operatorname{SL}_2(9)$ is $-I$: a determinant-one
involution in odd characteristic is diagonalizable with both eigenvalues
$-1$, unless it is the identity. Thus $f$ has order four and $f^2=-I$.
The polynomial $X^2+1$ has two distinct roots in $\F_9$, so $f$ fixes
two vector lines. On either line its scalar of order four generates
$\F_9$ over $\F_3$. It acts irreducibly on the additive two-dimensional
line. With line translations this gives a primitive line action;
flag transitivity transfers it to every line. Supergroups preserve both
local primitivity conditions.

Only the cyclic subgroup generated by $f$ is required to fix a direction.
The whole Sylow subgroup need not do so. This avoids the stronger
simultaneous Sylow-fixed-point assertion in the printed proof of
Proposition 3.9 while supplying its claimed conclusion.
\end{proof}

For the other Foulser plane cases, a faithful primitive action of $A_4$
or $S_4$ has degree at most four, by the normal Klein four subgroup.
The elementary maximal-subgroup list $A_4,D_{10},S_3$ for $A_5$ gives
primitive degrees $5,6,10$. Comparing the specified $q+1$ values in
source Proposition 3.7 leaves only $q=9$ among its exceptional cases.
